\documentclass[10pt,a4paper]{amsart}
\usepackage[margin=19mm]{geometry}
\usepackage{amsmath,amssymb,mathtools}
\usepackage[scr=rsfs,cal=euler]{mathalfa}
\usepackage{xfrac}
\usepackage[unicode,hidelinks,bookmarks=false]{hyperref}
\newcommand{\ie}{i.e.\ }
\newcommand{\cf}{cf.\ }
\newcommand{\resp}{respectively\ }
\newcommand{\Subsection}{\subsection}
\providecommand{\nobreakdash}{\nobreak\hbox{-}}
\setlength{\emergencystretch}{2em}
\multlinegap0pt
\usepackage{bm}
\usepackage{bbm}
\usepackage{dsfont}
\usepackage{xspace}
\usepackage{cancel}
\usepackage{mathtools}
\usepackage{amsthm}
\makeatletter
\@ifundefined{theorem}{\newtheorem{theorem}{Theorem}}{}
\makeatother
\title[The Dual-Server Domination Number of Some Graph Operations]{The Dual-Server Domination Number of Some Graph Operations}
\author{Shrilaxmi Laxminarayana Rao, Sayinath Udupa N.V.,, Prathviraj N}
\date{Source: arXiv:2608.14690v1 (CC BY 4.0)}
\begin{document}
\maketitle

\noindent\emph{Source and licence.} The Dual-Server Domination Number of Some Graph Operations. Version arXiv:2608.14690v1 (CC BY 4.0), distributed under the Creative Commons Attribution 4.0 International licence, as recorded in the arXiv licence metadata for this exact version and shown on the abstract page https://arxiv.org/abs/2608.14690v1. Full rights notice: licenses/arxiv-2608.14690-CC-BY-4.0.txt.\par\smallskip

\noindent\textit{Editorial scope.} Counted unit: the Theorem whose source label is \texttt{None} in the article (source0.tex lines 573--578, source label \texttt{None}), delivered together with the article's own complete original proof of it (source0.tex lines 580--705, one \texttt{proof} environment of 126 lines). No mathematics is written, repaired or re-derived: statement and proof are the article's own text line for line. Required context retained uncounted: none. Every definition, display and label of those lines is retained unchanged. Omitted from this package and not redistributed: 1--572, 579--579, 706--911. Nothing inside a retained excerpt is omitted or altered: every retained body is delivered exactly as the article's lines are written, so no omission receipt of any kind accompanies this package. Citations: the retained excerpts carry no citation command at all, so no bibliography entry of the article is transcribed and no reference is redistributed. Reference adaptation: none was needed and none was made; every reference inside the delivered unit resolves inside it, so no unresolved reference or citation is produced. Printed slips kept literal: no printed word, symbol, hypothesis or number of the counted statement or of its proof was altered. Layout and class: the article's own class is \texttt{article} with packages including graphicx, amsthm, hyperref, amsmath, amssymb, setspace, cite, ragged2e; the wrapper is the lane's pinned \texttt{amsart} editorial wrapper at 10pt on A4, which restates the theorem-like environments the retained text needs and adds the article's own macro definitions that the retained text needs (none), with extra wrapper packages (bm, bbm, dsfont, xspace, cancel, mathtools). The delivered numbering is the unit's own. Assets: no retained span contains a figure environment, a graphics-inclusion command, a PDF-page inclusion, a table or a TikZ picture; no image file is copied into this package, the package declares no asset dependency and no image file is redistributed with it. Licence: version arXiv:2608.14690v1 of this article is distributed under the Creative Commons Attribution 4.0 International licence, as recorded in the arXiv licence metadata for this exact version and shown on the abstract page https://arxiv.org/abs/2608.14690v1. A full rights notice is delivered at licenses/arxiv-2608.14690-CC-BY-4.0.txt. Characters: every retained excerpt is pure ASCII, so the delivered engine, XeLaTeX with its default Latin Modern text font, reports no missing character.
\begin{theorem}
Let $G$ and $H$ be graphs of orders $m\ge1$ and $n\ge1$, respectively. Then the dual-server domination number of the cartesian product  $G\times H $ is 
$\gamma_{ds}(G\times H)=mn$
if and only if $
\Delta(G)+\Delta(H)\le1$.
\end{theorem}

\begin{proof}
Let $
V(G)=\{u_1,u_2,\ldots,u_m\} \quad \text{ and }\quad
V(H)=\{v_1,v_2,\ldots,v_n\}$.
Then $
|V(G\times H)|=mn$.

Assume that $\gamma_{ds}(G\times H)=mn$. Let $S=R\cup B $
be a minimum dual-server dominating set of $G\times H$. Then $|S|=mn$. Now, suppose that  $\Delta(G)+\Delta(H)\ge2$.
We consider the following cases.

\textbf{Case 1.} $\Delta(G)\ge2$.\\
In particular, without loss of generality, let  $u_q$ and $u_r$ be two vertices adjacent to $u_k$. Choose any vertex $v_\ell\in V(H)$.

By the definition of the Cartesian product,
\[
(u_q,v_\ell)\sim(u_k,v_\ell)
\quad\text{and}\quad
(u_r,v_\ell)\sim(u_k,v_\ell).
\]
Hence the vertex $(u_k,v_\ell)$ has two distinct neighbours in
$G\times H$.
Without loss of generality, take
$
(u_q,v_\ell)\in R
\quad\text{and}\quad
(u_r,v_\ell)\in B.
$ Then $S \setminus (u_k,v_\ell) $  is not a DS-dominating set
, which implies
\[
|S| \neq mn,
\]
contradicting the assumption that $
\gamma_{ds}(G\times H)=mn $.

Therefore,
\[
\Delta(G)\le1.
\]

\textbf{Case 2.} $\Delta(H)\ge2$.\\
By an argument similar to that in Case~1, we obtain
\[
|S|\neq mn,
\]
contradicting the assumption that $
\gamma_{ds}(G\times H)=mn $.

Hence,
\[
\Delta(H)\le1.
\]

\textbf{Case 3.} $\Delta(G)=1$ and $\Delta(H)=1$.\\
Since $\Delta(G)=1$ and $\Delta(H)=1$, every vertex of $G$ and $H$ has at
most one neighbour.
Let $u_k$ be adjacent to $u_q$ in $G$, and let $v_\ell$ be adjacent to
$v_r$ in $H$. Then the four vertices $
(u_k,v_\ell),\;
(u_k,v_r),\;
(u_q,v_\ell),\;
(u_q,v_r)$
form a cycle of length $4$ in $G\times H$.\\
In particular, $(u_k,v_\ell)$ and $(u_q,v_r)$  adjacent to both
$(u_q,v_\ell)$ and $(u_k,v_r)$. Therefore, the vertex $(u_k,v_\ell)$  and $(u_q,v_r) $ have
two distinct neighbours in $G\times H$.\\
without loss of generality, take
\[
(u_q,v_\ell)\in R
\quad\text{and}\quad
(u_k,v_r)\in B.
\]
Hence, $
S\setminus\{(u_k,v_\ell),(u_q,v_r)\}
$ is not a dual-server dominating set. Hence 
\[
|S| \neq mn,
\]
which contradicts the assumption that
\[
\gamma_{ds}(G\times H)=mn.
\]

Thus the case $\Delta(G)=1$ and $\Delta(H)=1$ is impossible.

Therefore,
\[
\Delta(G)+\Delta(H)\le1.
\]

Conversely, suppose that 
$\Delta(G)+\Delta(H)\le1$.
Then one of the following cases must occur.

\textbf{Case 4.}  Suppose $\Delta(G)=0$ and $\Delta(H)=0$.\\
Then $V(G\times H)=\overline{K_{mn}} $. Therefore, \[
\gamma_{ds}(G\times H)=|V(G\times H)|=mn.
\]



\textbf{Case 5.} $\Delta(G)=1$ and $\Delta(H)=0$.\\
Since $\Delta(G)=1$, every vertex of $G$ has at most one neighbour.
Consequently, every vertex of $G\times H$ also has at most one neighbour.

Let $S=R\cup B$
be a DS-dominating set of $G\times H$. Every vertex outside $S$ must be adjacent to at
least one vertex in $R$ and $B$.

However, every vertex of $G\times H$ has degree at most one. Hence no
vertex can have two distinct neighbours. Therefore, no vertex can lie
outside $S$, and so $
S=V(G\times H)$.

Hence,
\[
\gamma_{ds}(G\times H)=|V(G\times H)|=mn.
\]

\textbf{Case 3.} $\Delta(G)=0$ and $\Delta(H)=1$.

By an argument similar to that in Case~2, we obtain
\[
\gamma_{ds}(G\times H)=mn.
\]
\end{proof}

\end{document}
